\subsection{Executable bindings and cache}

The current executor initially binds MemMode to the default compiled chain and ComputeMode to the native chain.
Both GEMM paths normalize the same FP16 product using FP32 RMSNorm statistics and weighting, then cast the final output to FP16.
Attention always binds to compiled SDPA plus RMSNorm.
A compiled-chain name does not establish that the matrix product and the complete normalization reduction occupy one kernel.

Calibration scopes distinguish analytical-only, ridge-only, the default ridge-plus-compute policy, and all-GEMM probing, retaining the $0.5R/2R$ analytical rule.
Cache entries use exact shape, strides, dtype, device, implementation identity, and device identity, alongside analytical bucket metadata.
An unseen signature must be prepared before regular timing or graph capture.
Every independently timed selector call invokes its actual chosen function.
Controlled graph replay binds static input storage and rejects calls with different storage.

\subsection{Complete-model integration}

The integration site is the actual attention output projection followed by residual addition and RMSNorm.
Attention, RoPE, the MLP, the KV cache, and the output head remain in the measured workload.
An explicit FP16 linear boundary preserves the matrix-product rounding position, and the compiled epilogue emulates precision casts.
HF and vLLM adapters are each validated against their own residual and normalization semantics; tolerances are not relaxed to mask a changed rounding contract.
Each policy prepares exact signatures before regular timing and graph capture, and is compared separately with the unmodified model and fixed-policy adapters.

\subsection{Correctness and timing obligations}

Candidates must pass complete-output checks with multiple random seeds and a larger-magnitude GEMM input before measurement.
The PyTorch model comparison additionally checks full logits on identical teacher-forced token histories.
Numerical rejection, unavailable execution, and valid timings are distinct recorded outcomes.
CUDA-event intervals and synchronized host batches are measured separately; preparation, calibration, and regular samples have separate records.

Section~\ref{sec:eval} specifies the main shapes, software versions, sample counts, and system workloads.
Legacy path names, copied selected medians, and split-region accounting are retained with the pilot protocol in Appendix~\ref{sec:pilot}.
Kernel names and launch counts are not substituted for DRAM bytes or classification accuracy.
